%! Simplifying Radicals; Adding \& Subtracting — Unit 5, Lesson 5.2 — Group Activity
\documentclass[10pt]{article}
\usepackage{algebra2-article}
\usepackage{algebra2-boxes}
\newcommand{\TallMath}[1]{$\displaystyle #1\rule[-1.4em]{0pt}{3.2em}$}

\begin{document}

\pageheader{Unit 5, Lesson 5.2}{Group Activity}
\namepartnerperiod

\vspace{0.08in}

\begin{headlinebox}{forestbg}
\small\textbf{Hiding in Plain Sight.} Two radicals that look nothing alike can be the same radical
wearing different clothes --- and you will never know until you simplify them. With your partner,
simplify \emph{first} and decide \emph{second}. Say out loud which perfect power you pulled out and
why it was the largest one. Assume all variables represent \textbf{positive} real numbers.
\end{headlinebox}

\vspace{0.06in}

\begin{tcolorbox}[colback=white, colframe=black!40, title=\textbf{Tier R --- Remediate},
    fonttitle=\bfseries, arc=1.5mm, left=3mm, right=3mm, top=2mm, bottom=2mm, breakable]
\textbf{1. Find the largest perfect square, then simplify.} The first row is done for you.
\begin{center}
\renewcommand{\arraystretch}{1.7}
\begin{tabularx}{\linewidth}{@{}l X X X@{}}
\hline
\textbf{Radical} & \textbf{Largest perfect square factor} & \textbf{Rewrite} & \textbf{Simplified} \\
\hline
$\sqrt{18}$  & $9$ & $\sqrt{9\cdot 2}$ & $3\sqrt2$ \\
$\sqrt{32}$  & \blank{1.6cm} & \blank{2.2cm} & \blank{2.0cm} \\
$\sqrt{44}$  & \blank{1.6cm} & \blank{2.2cm} & \blank{2.0cm} \\
$\sqrt{75}$  & \blank{1.6cm} & \blank{2.2cm} & \blank{2.0cm} \\
$\sqrt{90}$  & \blank{1.6cm} & \blank{2.2cm} & \blank{2.0cm} \\
$\sqrt{128}$ & \blank{1.6cm} & \blank{2.2cm} & \blank{2.0cm} \\
\hline
\end{tabularx}
\end{center}

\vspace{3pt}
\textbf{2. Finished, or not finished?} Mark each \textbf{D} (done --- already in simplest form) or
\textbf{N} (not done). If it is N, finish it.
\begin{center}
\renewcommand{\arraystretch}{1.7}
\begin{tabularx}{\linewidth}{@{}X X X@{}}
(a) \; $3\sqrt5$ \quad \blank{1.0cm} \; \blank{1.6cm} &
(b) \; $2\sqrt{12}$ \quad \blank{1.0cm} \; \blank{1.6cm} &
(c) \; $\sqrt{15}$ \quad \blank{1.0cm} \; \blank{1.6cm} \\
(d) \; $5\sqrt{8}$ \quad \blank{1.0cm} \; \blank{1.6cm} &
(e) \; $\sqrt[3]{9}$ \quad \blank{1.0cm} \; \blank{1.6cm} &
(f) \; $\sqrt[3]{24}$ \quad \blank{1.0cm} \; \blank{1.6cm} \\
\end{tabularx}
\end{center}

\vspace{3pt}
\textbf{3. Find the families.} Simplify all six, then sort them into two groups of like radicals.
\begin{center}
\renewcommand{\arraystretch}{1.6}
\begin{tabularx}{\linewidth}{@{}X X X@{}}
$\sqrt{12}=$ \blank{1.8cm} & $\sqrt{50}=$ \blank{1.8cm} & $\sqrt{27}=$ \blank{1.8cm} \\
$\sqrt{8}=$ \blank{1.8cm} & $\sqrt{75}=$ \blank{1.8cm} & $\sqrt{32}=$ \blank{1.8cm} \\
\end{tabularx}
\end{center}
\begin{itemize}[leftmargin=1.6em, itemsep=3pt]
  \item The $\sqrt{3}$ family: \blank{5.5cm} \; Their sum: \blank{2.0cm}
  \item The $\sqrt{2}$ family: \blank{5.5cm} \; Their sum: \blank{2.0cm}
\end{itemize}

\vspace{2pt}
\textbf{4.} Why can the two sums above not be combined into a single term? \writeline
\end{tcolorbox}

\begin{tcolorbox}[colback=white, colframe=black!40, title=\textbf{Tier A --- Approaching Proficiency},
    fonttitle=\bfseries, arc=1.5mm, left=3mm, right=3mm, top=2mm, bottom=2mm, breakable]
\textbf{1. Simplify completely.}
\begin{center}
\renewcommand{\arraystretch}{2.0}
\begin{tabularx}{\linewidth}{@{}X X X@{}}
(a) \; $\sqrt{162}=$ \blank{2.0cm} &
(b) \; $\sqrt[3]{128}=$ \blank{2.0cm} &
(c) \; $\sqrt{48x^{7}}=$ \blank{2.2cm} \\
(d) \; $\sqrt[3]{81y^{8}}=$ \blank{2.2cm} &
(e) \; $\sqrt{\dfrac{45}{4}}=$ \blank{2.0cm} &
(f) \; $\dfrac{\sqrt{98}}{\sqrt{2}}=$ \blank{2.0cm} \\
\end{tabularx}
\end{center}

\vspace{3pt}
\textbf{2. Combine.} Simplify every term \emph{first}, then collect like radicals.
\begin{center}
\renewcommand{\arraystretch}{2.0}
\begin{tabularx}{\linewidth}{@{}X X@{}}
(a) \; $2\sqrt{18}+3\sqrt{50}=$ \blank{2.6cm} &
(b) \; $\sqrt{80}-\sqrt{45}+\sqrt{20}=$ \blank{2.6cm} \\
(c) \; $\sqrt{27x}+\sqrt{12x}=$ \blank{2.6cm} &
(d) \; $4\sqrt{24}-\sqrt{54}+\sqrt{6}=$ \blank{2.6cm} \\
\end{tabularx}
\end{center}

\vspace{3pt}
\textbf{3. Find the mistake.} Each student made exactly one error. Name it and give the correct
answer.
\begin{center}
\renewcommand{\arraystretch}{1.7}
\begin{tabularx}{\linewidth}{@{}l X X@{}}
\hline
\textbf{Student work} & \textbf{What went wrong} & \textbf{Correct answer} \\
\hline
$\sqrt{20}+\sqrt{5}=\sqrt{25}=5$ & \blank{4.2cm} & \blank{2.2cm} \\
$3\sqrt7+2\sqrt7=5\sqrt{14}$ & \blank{4.2cm} & \blank{2.2cm} \\
$\sqrt[3]{54}=\sqrt[3]{9\cdot 6}=3\sqrt[3]{6}$ & \blank{4.2cm} & \blank{2.2cm} \\
\hline
\end{tabularx}
\end{center}

\vspace{3pt}
\textbf{4. Exact beats approximate.} A rectangular garden plot measures $\sqrt{75}$ meters by
$\sqrt{12}$ meters.
\begin{enumerate}[label=(\alph*), leftmargin=1.9em, itemsep=4pt]
  \item Write each side in simplest radical form: \; length \blank{2.0cm}, \; width \blank{2.0cm}
  \item Find the \textbf{exact} perimeter as a single simplified radical: \blank{2.6cm}
  \item Now approximate it to the nearest tenth of a meter: \blank{2.4cm}
  \item Why would a builder want the answer in part (b) and not just the one in part (c)?
        \writeline
\end{enumerate}
\end{tcolorbox}

\begin{tcolorbox}[colback=white, colframe=black!40, title=\textbf{Tier E --- Extension},
    fonttitle=\bfseries, arc=1.5mm, left=3mm, right=3mm, top=2mm, bottom=2mm, breakable]
\textbf{Part 1 --- Prove the error is never a lucky guess.} Everyone knows
$\sqrt{a}+\sqrt{b}\ne\sqrt{a+b}$ ``usually.'' You are going to prove it is \emph{never} true unless
one of the two numbers is zero. Let $a\ge 0$ and $b\ge 0$, and \emph{suppose} for a moment that
\[
  \sqrt{a}+\sqrt{b}=\sqrt{a+b}.
\]
\begin{enumerate}[label=(\alph*), leftmargin=1.9em, itemsep=5pt]
  \item Square both sides. The left side is a binomial squared --- expand it carefully.
        \writeline
  \item Subtract $a+b$ from both sides. What is left? \blank{3.4cm}
  \item Solve that equation. What must be true about $a$ and $b$? \blank{3.6cm}
  \item State the conclusion in a sentence, then check it against $a=9$, $b=16$. \writelines{2}
\end{enumerate}

\vspace{5pt}
\textbf{Part 2 --- A staircase of radicals.} Simplify each term, then add.
\[
  \sqrt{2}+\sqrt{8}+\sqrt{18}+\sqrt{32}+\sqrt{50} \;=\;
  \blank{1.2cm}+\blank{1.2cm}+\blank{1.2cm}+\blank{1.2cm}+\blank{1.2cm}
  \;=\; \blank{2.0cm}
\]
\begin{enumerate}[label=(\alph*), leftmargin=1.9em, itemsep=5pt]
  \item Look at the radicands $2,\,8,\,18,\,32,\,50$. Write each one in the form
        $2\cdot(\text{something})^{2}$: \blank{6.0cm}
  \item Explain in one sentence why $\sqrt{2k^{2}}=k\sqrt{2}$. \writeline
  \item What is the \emph{next} term in the staircase, and what does it simplify to?
        \blank{3.6cm}
  \item Predict the total of the first $n$ terms, in terms of $n$: \blank{3.0cm} \\
        \emph{Hint:} you added $1+2+3+\dots$ Check your formula against the $n=5$ answer above.
        \writeline
  \item Build the same staircase with $3$ in place of $2$: simplify
        $\sqrt{3}+\sqrt{12}+\sqrt{27}+\sqrt{48}=\blank{2.4cm}$
\end{enumerate}

\vspace{5pt}
\textbf{Part 3 --- Two functions that were secretly the same shape.} Consider
$f(x)=\sqrt{8x}$ and $g(x)=\sqrt{2x}$.
\begin{center}
\renewcommand{\arraystretch}{1.3}
\begin{tabular}{|c|c|c|c|}
\hline
$x$ & $2$ & $8$ & $18$ \\ \hline
$g(x)=\sqrt{2x}$ & \blank{1.2cm} & \blank{1.2cm} & \blank{1.2cm} \\ \hline
$f(x)=\sqrt{8x}$ & \blank{1.2cm} & \blank{1.2cm} & \blank{1.2cm} \\ \hline
\end{tabular}
\end{center}
\begin{enumerate}[label=(\alph*), leftmargin=1.9em, itemsep=5pt]
  \item In every column, $f(x)$ is exactly \blank{1.6cm} times $g(x)$.
  \item Prove it with the product property: $\sqrt{8x}=\sqrt{\blank{1.0cm}\cdot 2x}=\blank{2.0cm}$
  \item So $f$ is not a new function at all --- it is $g$ stretched \textbf{vertically} by a factor
        of \blank{1.0cm}. Which of the two do $f$ and $g$ share: the same domain, the same
        endpoint, or both? \writeline
\end{enumerate}
\end{tcolorbox}

\end{document}
